%%%%%%%%%%%%%%%%%%%%%%%%%%%%%%%%%%%%%%%%%
% FRI Data Science_report LaTeX Template
% Version 1.0 (28/1/2020)
% 
% Jure Demšar (jure.demsar@fri.uni-lj.si)
%
% Based on MicromouseSymp article template by:
% Mathias Legrand (legrand.mathias@gmail.com) 
% With extensive modifications by:
% Antonio Valente (antonio.luis.valente@gmail.com)
%
% License:
% CC BY-NC-SA 3.0 (http://creativecommons.org/licenses/by-nc-sa/3.0/)
%
%%%%%%%%%%%%%%%%%%%%%%%%%%%%%%%%%%%%%%%%%


%----------------------------------------------------------------------------------------
%	PACKAGES AND OTHER DOCUMENT CONFIGURATIONS
%----------------------------------------------------------------------------------------
\documentclass[fleqn,moreauthors,10pt]{ds_report}
\usepackage[english]{babel}

\graphicspath{{fig/}}




%----------------------------------------------------------------------------------------
%	ARTICLE INFORMATION
%----------------------------------------------------------------------------------------

% Header
\JournalInfo{FRI Natural language processing course 2026}

% Interim or final report
\Archive{Project report} 
%\Archive{Final report} 

% Article title
\PaperTitle{Translation assistant for legal terminology} 

% Authors (student competitors) and their info
\Authors{Romina Mihalič, Anakhita Khademipur, Maša Pipan, Luka Škoda}

% Advisors
\affiliation{\textit{Advisors: Aleš Žagar}}

% Keywords
\Keywords{translation, Slovenian, English, terminology}
\newcommand{\keywordname}{Keywords}


%----------------------------------------------------------------------------------------
%	ABSTRACT
%----------------------------------------------------------------------------------------

\Abstract{

	For this project we are going to create an AI tranlsation assistant, that will be able to translate between Slovenian and English. We will use the T5 transformer architecture, which is a state-of-the-art model for natural language processing tasks. We will train the model on a large dataset of parallel sentences in Slovenian and English, and we will evaluate its performance on a test set of sentences that it has not seen before. We will also compare our model's performance to other translation models, such as Google Translate and DeepL. The model focuses on legal terminology.

}

%----------------------------------------------------------------------------------------

\begin{document}

% Makes all text pages the same height
\flushbottom 

% Print the title and abstract box
\maketitle 

% Removes page numbering from the first page
\thispagestyle{empty} 

%----------------------------------------------------------------------------------------
%	ARTICLE CONTENTS
%----------------------------------------------------------------------------------------

\section*{Introduction}

Accurate translation in specialized domains such as legal and administrative texts is difficult, especially when using general machine translation tools. While systems like Google Translate or large language models can produce fluent text, they often make mistakes with domain-specific terminology or use inconsistent translations. This is not ideal because in institutional contexts precise and standardized terminology is key for clarity and correctness. 

Since the introduction of the Transformer architecture, Transformer models have revolutionized the field of machine translation and have replaced older models, such as CNNs and RNNs that process data word by word, unlike Transformer models that work with the text as a whole. Some recently released large language models such as Llama and Gemma adapt this concept and provide more flexible and context-aware text generation, nonetheless domain-specific translation is still difficult because of a lack of consistent terminology and factual accuracy. 

In recent years, new approaches such as Retrieval-Augmented Generation (RAG) have been proposed to improve this. RAG combines information retrieval with text generation: instead of translating only based on the model, it first retrieves relevant documents and then uses them to produce a more accurate and grounded output.  

There are also many existing tools and frameworks for language processing, such as NLTK or spaCy, but they are general-purpose and do not focus on domain-specific translation.  

For this project, we focus on the institutional and administrative domain and translation from English into Slovenian. The main corpus used is the DGT Translation Memory, which contains officially translated documents from the European Commission. These texts include high-quality legal, administrative and policy-related parallel data in 24 official EU languages, including English and Slovenian, and are suitable for extracting domain-specific terminology and context and improving translation quality.  

The goal of this project is to build a translation assistant that produces more accurate and consistent legal translations based on retrieved institutional texts. 
%------------------------------------------------

\section*{Methods}

First of our proposed corpora is the DGT Translation Memory, which is a multilingual parallel corpus made out of .tmx documents from the European Commission Agency\cite{DGT_Translation_Memory}. For this project we will extract only the English and Slovenian segments. We will supplement this with aligned English-Slovenian texts from EUR-Lex, which includes EU treaties, regulations and directives\cite{EUR-Lex}. For terminology reference we will use the IATE - a bilingual glossary of legal and administrative terms\cite{IATE_-_Interactive_Terminology_for_Europe}. 

\section*{Implementation}
We constructed our dataset from the DGT Translation Memory, which contains parallel texts from European Union institutions. From this corpus, we extracted aligned English–Slovenian sentence pairs and converted them into a structured CSV file format for easier processing.
We randomly sampled a subset of 500 sentence pairs, which were used for translation and evaluation during this phase. Additionally, we created a separate test subset of 200 sentence pairs, which will be used later on unseen data.

The dataset consists of parallel sentences, where each English sentence is paired with its corresponding Slovenian translation. The Slovenian sentences serve as the reference (gold standard) for evaluation.

We implemented a baseline machine translation system using a pretrained transformer model. Instead of training a model from scratch, we used the NLLB-200 model (facebook/nllb-200-distilled-600M), which supports multilingual translation including English and Slovenian.

The model and tokenizer were loaded using the HuggingFace Transformers library. We configured the source language as English (eng\_Latn) and the target language as Slovenian (slv\_Latn) to ensure correct translation direction.

To improve performance and reduce runtime, translations were processed in batches of eight sentences at a time. This batching strategy significantly speeds up inference compared to processing each sentence individually.

\section*{Evaluation}

For quantitative evaluation, we used the BLEU (Bilingual Evaluation Understudy) metric, implemented via the SacreBLEU library. BLEU measures the similarity between machine-generated translations and reference translations based on overlapping n-grams.

The model achieved a BLEU score of 36.17 on the sampled dataset. This result indicates that the model produces reasonably accurate translations, though there is still room for improvement, particularly in handling specialized legal terminology.
%----------------------------------------------------------------------------------------
%	REFERENCE LIST
%----------------------------------------------------------------------------------------
\bibliographystyle{unsrt}
\bibliography{report}


\end{document}